% ===================================================================
%  ТАБЛИЦА № 4 — Зрелость и Старость.
%  Traduction 2026 d'apres texto/io, controlee sur texto/fr, texto/en
%  et texto/es. Consigne : voir texto/ru/10-tablica-01.tex.
% ===================================================================

%%K t04-apar-1 apar
\VUsaut{4.0mm}
\VUtitre{15.0pt}{60}{ТАБЛИЦА № 4}
\VUsaut{1.6mm}
\VUtitre{11.4pt}{0}{Зрелость и Старость.}
\VUsaut{3.2mm}

%%K t04-tit-1 sub
\VUcentre{11.4pt}{0}{\textit{Первая сцена.}}
\VUsaut{1.6mm}
\VUtitre{11.4pt}{0}{Свадебный обед.}
\VUsaut{2.6mm}

%%K t04-tit-2 sub
\VUpk{10.2pt}{40}{Осень.}
\VUsaut{3.0mm}

%%K t04-c1-01-1 p
1. --- Молодые люди выросли. Один из них, Иван, женится. Мы на
\VUgras{свадебном обеде}, который собрал за одним столом семьи жениха и
невесты.

%%K t04-c1-02-1 p
2. --- Стоит \VUgras{осень}, месяц \VUgras{сентябрь}. Холодный ветер
\VUgras{октября} и \VUgras{ноября} ещё не сорвал с полей зелёной листвы.
Вечерний воздух мягок и душист; поэтому обед подают под \VUgras{верандой}
\textsuperscript{(8)}, выходящей в сад.

%%K t04-c1-03-1 p
3. --- Вдали, на \VUgras{холме} \textsuperscript{(3)}, стоит дом Ивановых
родителей, изящная \VUgras{усадьба} \textsuperscript{(1)}, укрытая в зелени;
прекрасные виноградники, дающие отличное вино, покрывают склон холма.
Виноградарь ходит за ними и холит их.

%%K t04-c1-04-1 p
4. --- Над лозами пролетает стайка \VUgras{куропаток} \textsuperscript{(2)},
вспугнутая приближением \VUgras{лесника} \textsuperscript{(5)}. А тот, с
\VUgras{ружьём} под мышкой и \VUgras{ягдташем через плечо}, обшаривает
\VUgras{кусты} \textsuperscript{(4)} со своей \VUgras{собакой} \textsuperscript{(6)}.
Он стережёт поместье и не даёт браконьерам истреблять дичь.

%%K t04-c1-05-1 p
5. --- Снаружи по \VUgras{веранде} вьётся \VUgras{шпалерная лоза}, с которой
свисают тяжёлые \VUgras{грозди} \textsuperscript{(7)}.

%%K t04-c1-06-1 p
6. --- Прекрасные осенние цветы, украшающие \VUgras{стол} \textsuperscript{(16)},
--- это \VUgras{хризантемы} \textsuperscript{(9)}; \VUgras{отец} \textsuperscript{(10)}
невесты вдел одну в петлицу фрака.

%%K t04-c1-07-1 p
7. --- Обед подходит к концу. \VUgras{Слуга} \textsuperscript{(11)} начинает
вносить блюда десерта и скоро уберёт ставшие ненужными части приборов ---
\VUgras{ложки} \textsuperscript{(14)}, \VUgras{вилки} \textsuperscript{(12)},
\VUgras{ножи} \textsuperscript{(13)}, \VUgras{блюда} \textsuperscript{(15)}.

%%K t04-tit-3 sub
\VUpk{10.2pt}{40}{Семья.}
\VUsaut{3.0mm}

%%K t04-c1-08-1 p
8. --- У всех вид весёлый, а улыбка, с какой \VUgras{мать} \textsuperscript{(17)}
жениха глядит на своих детей, показывает, как она счастлива.

%%K t04-c1-09-1 p
9. --- Этот брак соединяет две богатые семьи округи. Иван, \VUgras{муж}
\textsuperscript{(18)}, --- \VUgras{сын} крупного землевладельца, и он становится
\VUgras{зятем} преуспевающего фабриканта.

%%K t04-c1-10-1 p
10. --- \VUgras{Супруги} молоды, и всё же они долго были
\VUgras{помолвлены}. \VUgras{Молодая жена} \textsuperscript{(19)}, Марта, мила в
своём \VUgras{белом платье}, убранном флёрдоранжем.

%%K t04-c1-11-1 p
11. --- Её \VUgras{свёкор} \textsuperscript{(20)} рад такой доброй
\VUgras{невестке}, а \VUgras{тёща} \textsuperscript{(21)} Ивана улыбается, видя
свою \VUgras{дочь} такой красивой.

%%K t04-c1-12-1 p
12. --- На конце стола сидят прочие \VUgras{гости} \textsuperscript{(22)}:
\VUgras{родня, друзья, двоюродные}. \VUgras{Метрдотель} \textsuperscript{(23)} с
салфеткой на руке проворно им прислуживает.

%%K t04-tit-4 sub
\VUpk{10.2pt}{40}{Друзья.}
\VUsaut{3.0mm}

%%K t04-c1-13-1 p
13. --- По правую сторону стола --- \VUgras{барышня} \textsuperscript{(24)},
\VUgras{подруга} невесты, затем \VUgras{брат} невесты, её
\VUgras{шафер} \textsuperscript{(25)}. Он беседует со своей \VUgras{подружкой
невесты} \textsuperscript{(26)}, сестрой жениха, которая через этот брак
становится \VUgras{золовкой} Марты. Это прелестная молодая женщина. У неё
красивые \VUgras{драгоценности}, \VUgras{жемчужное ожерелье} и золотой
\VUgras{браслет}.

%%K t04-c1-14-1 p
14. --- Рядом с ними господин, который встал и поднял бокал, ---
\VUgras{друг} \textsuperscript{(27)} семьи. Он один из \VUgras{свидетелей жениха}.
Он говорит \VUgras{тост}, пьёт за \VUgras{здоровье} молодых и желает им
долгого счастья. Он в парадном платье: фрак и \VUgras{лакированные туфли}
\textsuperscript{(28)}. Он ведёт под руку \VUgras{бабушку} \textsuperscript{(29)},
которая \VUgras{вдова}.

%%K t04-c1-15-1 p
15. --- Молодой человек во \VUgras{фраке} \textsuperscript{(31)}, с тонкими
чёрными усами, --- \VUgras{шурин} \textsuperscript{(30)} Ивана. Он видный
инженер. Он сидит рядом с очень изящной \VUgras{дамой} \textsuperscript{(33)},
\VUgras{старшей сестрой} Марты и матерью той милой девочки, которая идёт
под руку с двоюродным братом поднести цветок и \VUgras{пожелать}
\VUgras{доброго вечера}. У этой дамы прекрасные \VUgras{серьги} и изящная
\VUgras{причёска}.

%%K t04-c1-16-1 p
16. --- Маленький двоюродный брат, такой смешной в своей \VUgras{блузе}
\textsuperscript{(36)} и широких \VUgras{штанишках} \textsuperscript{(37)}, очень
несчастный. Он \VUgras{сирота} \textsuperscript{(35)}. Он потерял отца и
мать совсем малым. Его дядя --- его \VUgras{опекун} \textsuperscript{(38)} и
остался холостым, чтобы вырастить бедного ребёнка. Он добрый человек. Он
изрядно лыс и очень близорук; он носит \VUgras{пенсне}.

%%K t04-tit-5 sub
\VUsaut{2.4mm}
\VUpk{10.2pt}{40}{Десерт.}
\VUsaut{3.0mm}

%%K t04-c1-17-1 p
17. --- Позади гостей, на столе, покрытом \VUgras{скатертью}, разложен
\VUgras{десерт} \textsuperscript{(39)}. В большой вазе --- фрукты:
\VUgras{персики} \textsuperscript{(40)}, \VUgras{груши} \textsuperscript{(41)},
\VUgras{яблоки} \textsuperscript{(42)}, \VUgras{абрикосы} \textsuperscript{(43)} и
\VUgras{виноград} \textsuperscript{(44)}. Рядом --- \VUgras{ананасы}
\textsuperscript{(45)}, прекрасный \VUgras{торт} \textsuperscript{(46)},
\VUgras{бутылки ликёра} \textsuperscript{(48)} и \VUgras{шампанского}
\textsuperscript{(49)}, а также стопки чистых \VUgras{тарелок} \textsuperscript{(47)}.

%%K t04-c1-18-1 p
18. --- \VUgras{Виночерпий} \textsuperscript{(50)} откупоривает \VUgras{бутылку}
\textsuperscript{(52)} старого вина \VUgras{штопором} \textsuperscript{(51)}.
\VUgras{Пробка} \textsuperscript{(53)} идёт, кажется, очень туго. У этого
виночерпия на руке \VUgras{салфетка} \textsuperscript{(54)}.

%%K t04-c1-19-1 p
19. --- После обеда мужчины перейдут в курительную, а дамы пойдут
одеваться к балу.

%%K t04-tit-6 sub
\VUsaut{5.0mm}
\VUcentre{11.4pt}{0}{\textit{Вторая сцена.}}
\VUsaut{1.6mm}
\VUpk{11.4pt}{40}{Именины дедушки.}
\VUsaut{2.6mm}

%%K t04-tit-7 sub
\VUpk{10.2pt}{40}{Визит.}
\VUsaut{3.0mm}

%%K t04-c2-01-1 p
1. --- Прошло десять лет. Родители Ивана состарились и нежно любят своих
троих внуков: двух \VUgras{внуков} \textsuperscript{(2)} и \VUgras{внучку}
\textsuperscript{(3)}. Так как сегодня \VUgras{дедушкины именины}, дети пришли
его поздравить.

%%K t04-c2-02-1 p
2. --- Маленький Пётр ещё совсем мал; ему только три года. Он видит, как
к нему идёт \VUgras{кошка} \textsuperscript{(1)}. Ему хочется погладить её
красивую шелковистую \VUgras{шерсть} и длинный \VUgras{хвост}; ему и в
голову не приходит, что под этими изящными \VUgras{лапками} есть острые
злые коготки, которые могут его оцарапать.

%%K t04-c2-03-1 p
3. --- Его сестра Габриэль, которая держит его за руку, куда серьёзнее, а
Марсель в большом \VUgras{пальто} \textsuperscript{(4)} и кожаном \VUgras{поясе}
\textsuperscript{(5)} выглядит совсем взрослым, несмотря на короткие штаны, из
которых видны его \VUgras{чулки} \textsuperscript{(6)}.

%%K t04-c2-04-1 p
4. --- Их отец надел толстое зимнее \VUgras{пальто} \textsuperscript{(7)} с
\VUgras{меховым воротником} \textsuperscript{(8)}, а в \VUgras{кармане}
\textsuperscript{(9)} у него шарф. На его жене войлочная шляпа с \VUgras{перьями}
\textsuperscript{(10)}, \VUgras{жакет} \textsuperscript{(11)}, меховое \VUgras{боа}
\textsuperscript{(12)} и \VUgras{муфта} \textsuperscript{(13)}.

%%K t04-c2-05-1 p
5. --- Очень холодно: пришла зима. \VUgras{Снег} \textsuperscript{(19)} шёл
весь день, и крыши домов совсем белы. С \VUgras{ночью} мороз крепчает ещё
пуще. В небе светят \VUgras{луна} \textsuperscript{(17)} и \VUgras{звёзды}
\textsuperscript{(18)}, пронизывая \VUgras{тьму}.

%%K t04-tit-8 sub
\VUsaut{4.2mm}
\VUpk{10.2pt}{40}{Болезнь.}
\VUsaut{3.0mm}

%%K t04-c2-06-1 p
6. --- Бедный \VUgras{слепой} \textsuperscript{(20)}, который переходит площадь
перед \VUgras{аптекой} \textsuperscript{(16)}, ведомый своим \VUgras{пуделем}
\textsuperscript{(21)}, очень несчастен. Жюстина, \VUgras{служанка}
\textsuperscript{(14)}, несёт из кухни \VUgras{чашку} \textsuperscript{(15)} очень
горячего \VUgras{травяного настоя}, щедро подслащённого.

%%K t04-c2-07-1 p
7. --- \VUgras{Бабушка} \textsuperscript{(23)}, которой хочется взять со стола
свои \VUgras{очки} \textsuperscript{(26)}, чтобы прочесть \VUgras{рецепт}
\textsuperscript{(27)}, только что написанный \VUgras{врачом} \textsuperscript{(25)},
поднимается из кресла, опираясь на свою \VUgras{палку} \textsuperscript{(24)}.

%%K t04-c2-08-1 p
8. --- Она часто страдает мелкими \VUgras{недугами}, \VUgras{головокружениями},
\VUgras{простудами} и \VUgras{зубной болью}; ноги у неё слабы, а глаза уже
неважны. И всё же она рада подразнить своего старого \VUgras{доктора},
который всегда хочет прописать ей \VUgras{микстуру} \textsuperscript{(28)}.
Сегодня она очень счастлива, что молодая родня вокруг неё.

%%K t04-c2-09-1 p
9. --- В хорошо закрытой комнате уютно. В мраморном \VUgras{камине}
\textsuperscript{(29)} пылает \VUgras{славный огонь} \textsuperscript{(30)}, и хорошо
на душе при мягком \VUgras{свете} только что зажжённой \VUgras{лампы}
\textsuperscript{(31)}.

%%K t04-tit-9 sub
\VUsaut{4.2mm}
\VUpk{10.2pt}{40}{Дедушка.}
\VUsaut{3.0mm}

%%K t04-c2-10-1 p
10. --- Даже \VUgras{дедушка} \textsuperscript{(33)} забыл, что был болен, что
\VUgras{ревматизм} держит его прикованным к \VUgras{креслу} \textsuperscript{(34)}
и что ещё вчера у него были \VUgras{жар и обмороки}. Он уже не помнит, что
прошлой зимой у него было \VUgras{воспаление лёгких} и что жестокий,
мучительный \VUgras{кашель} заставил его немало настрадаться.

%%K t04-c2-10-2 p
А ведь оправился он с трудом. Теперь ему получше. Но нога, которую он
держит на \VUgras{подушке} \textsuperscript{(38)}, положенной на низкую
\VUgras{скамеечку} \textsuperscript{(39)}, всё ещё болит у него.

%%K t04-c2-11-1 p
11. --- Когда он заботливо укутан в свой тёплый \VUgras{халат}
\textsuperscript{(35)} с большими перламутровыми \VUgras{пуговицами}
\textsuperscript{(36)} и когда рядом с ним, чтобы читать ему газету, сидит
маленькая \VUgras{сиротка} \textsuperscript{(40)} в белом \VUgras{чепчике}, в
синем \VUgras{платье} \textsuperscript{(42)} и \VUgras{пелеринке}
\textsuperscript{(41)}, он не жалуется.

%%K t04-c2-12-1 p
12. --- Каждый год, когда \VUgras{календарь} \textsuperscript{(43)} снова
приводит \VUgras{число} первого \VUgras{декабря}, он с нетерпением ждёт
своих \VUgras{внуков} --- ведь он знает, что они не забудут этого важного
\VUgras{числа}.

%%K t04-c2-13-1 p
13. --- Он с чувством вспоминает далёкие дни, когда сам ходил
поздравлять славного императорского \VUgras{маршала}, чей \VUgras{портрет}
\textsuperscript{(44)} висит на стене и который приходится детям
\VUgras{прадедом}.
\VUsaut{6.0mm}
\VUfilet{22mm}
